\subsection{Data and Excess Growth Rates}
\label{sec:data}

The single-asset analysis uses daily SPDR S\&P~500 ETF (SPY) prices from January~3, 2014 to January~3, 2024 (a ten-year training window, $T_{\text{IS}} = 2{,}516$), with the period from January~4, 2024 through April~20, 2026 ($T_{\text{OoS}} = 572$) held out. Prices are split-adjusted but not dividend-adjusted; the data sources are described in the acknowledgments. For ticker $i$ and trading day $j$ the annualised excess growth rate is
\begin{equation}
    G_{i,j} \equiv \frac{1}{\Delta t} \ln\!\left(\frac{P_{i,j}}{P_{i,j-1}}\right) - r_f, \qquad \Delta t = 1/252, \quad r_f = 0,
    \label{eq:growth_rate}
\end{equation}
with $P_{i,j}$ the session volume-weighted average price (VWAP). The factor $1/\Delta t$ is a constant rescaling; all simulations and metrics are evaluated on the daily-scale series. The empirical study is organised along two pipelines (full schematic in Appendix~\ref{sec:supp_pipeline_schematic}). \emph{Pipeline~A} fits a CHMM to each ticker independently and evaluates single-asset metrics; \emph{Pipeline~B} reuses the per-asset Pipeline-A fits as marginals and injects cross-asset dependence through a rank-based Student-$t$ copula (\S\ref{sec:cross_asset_methods}).

\subsection{Continuous Hidden Markov Model}
\label{sec:chmm_def}

We model $G_t$ as a continuous HMM $\mathcal{M} = (\mathcal{S}, \mathbf{T}, \mathbf{B}, \boldsymbol{\pi})$ with $K$ latent states, per-state emission density $b_k(x; \boldsymbol\theta_k)$, transition matrix $\mathbf{T} \in \mathbb{R}^{K \times K}$, and initial distribution $\boldsymbol{\pi}$ \cite{rabiner1986introduction}. The stationary marginal is the $K$-component mixture
\begin{equation}
    f(x) = \sum_{k=1}^{K} \bar{\pi}_k\, b_k(x; \boldsymbol\theta_k), \qquad \bar{\boldsymbol{\pi}} = \bar{\boldsymbol{\pi}} \mathbf{T},
    \label{eq:mixture}
\end{equation}
with $\bar{\boldsymbol\pi}$ the stationary distribution. Existence and uniqueness require irreducibility and aperiodicity of $\mathbf T$ (Assumption~\ref{ass:irred} in \S\ref{sec:theory}). We learn $\mathbf{T}$, $\boldsymbol\theta_{1:K}$, and $\boldsymbol{\pi}$ jointly by EM (\S\ref{sec:estimation}).

The four emission families share a single forward-backward scaffold and quantile-based initialization; the M-step is the only architectural difference. \emph{CHMM-N}: $b_k = \mathcal{N}(\mu_k, \sigma_k^2)$. \emph{CHMM-t}: $b_k = t_{\nu_k}(\mu_k, \sigma_k)$ with per-state degrees of freedom $\nu_k \in [\nu_{\min}, \nu_{\max}]$. \emph{CHMM-L}: $b_k = \mathrm{Laplace}(\mu_k, b_k)$. \emph{CHMM-GED}: $b_k = \mathrm{GED}(\mu_k, \alpha_k, p_k)$ with per-state shape $p_k \in [p_{\min}, p_{\max}]$ that nests Gaussian ($p_k = 2$) and Laplace ($p_k = 1$). The architectural diagram is in Appendix~\ref{sec:supp_architecture}.

\subsection{Benchmark Generators}
\label{sec:benchmarks}

We benchmark against representatives from each generator tier, all fit on the same IS window: i.i.d.\ bootstrap \cite{politis1994stationary} (non-parametric distributional ceiling), Gaussian and Laplace i.i.d., GARCH(1,1) Gaussian \cite{bollerslev1986generalized} and Student-$t$ \cite{bollerslev1987conditionally}, and Markov-switching GARCH (MS-GARCH; \citealp{haas2004new}). The extended GARCH-family panel (EGARCH, GJR-GARCH, HAR-RV, MS-GARCH at $K \in \{3, 4, 6\}$), an explicit-duration semi-Markov foil, a deep-generative QuantGAN \cite{wiese2020quantgan} row, plus stochastic-volatility (Taylor-Harvey-Ruiz-Shephard), Markov-switching multifractal \cite{calvet2004regime}, and Merton-jump-diffusion \cite{merton1976option} baselines are reported in Appendix~\ref{sec:extended_baselines} and \ref{sec:sv_msm_jd_baselines}.

\subsection{Cross-Asset Dependence (Pipeline~B)}
\label{sec:cross_asset_methods}

Pipeline~B reuses each per-asset CHMM marginal and injects cross-asset dependence through a rank-based Student-$t$ copula \cite{demarta2005tcopula, mcneil2015quantitative}. Pseudo-uniform observations $u_{j,t} = \mathrm{rank}(g_{j,t}) / (T+1)$ feed an analytic Kendall's-$\tau$ correlation estimate $\rho_{ij} = \sin(\pi \tau_{ij}/2)$ \cite{embrechts2002correlation}; profile MLE selects $\nu^\star$. Sklar's theorem \cite{sklar1959fonctions} guarantees that the rank-reordering simulator of \citet{iman1982distribution} preserves each fitted CHMM marginal exactly while coupling the asset ranks to the copula sample. SIM, Gaussian-copula, and truncated C-vine comparators are in Appendix~\ref{sec:cross_asset_supp}.

\subsection{Evaluation Protocol}
\label{sec:metrics}

For each model we simulate $P = 1{,}000$ independent paths of length $T$ following \cite{stenger2024thinking}, under a deterministic global seed root $20260420$. Headline metrics are the two-sample Kolmogorov--Smirnov \cite{kolmogorov1933, smirnov1948} pass rate at $\alpha = 0.05$, mean simulated excess kurtosis, and the absolute-return ACF mean absolute error (MAE) over $L = 252$ lags,
\begin{equation}
    \text{ACF-MAE} = \frac{1}{L}\sum_{\tau=1}^{L} \left| \hat\rho^{\text{obs}}_{|G|}(\tau) - \bar{\hat\rho}^{\,\text{sim}}_{|G|}(\tau) \right|.
    \label{eq:acf_mae}
\end{equation}
A separate Value-at-Risk panel (\S\ref{sec:var_backtest}) carries the unconditional Kupiec \cite{kupiec1995techniques} and Expected-Shortfall envelope diagnostics. Block-bootstrap KS recalibration, Anderson--Darling, $W_1$, Hellinger, and CRPS auxiliaries are in Appendix~\ref{sec:supp_metrics}.
